\section{Underlying Functions and Addresses}
This section defines the cryptographic primitives and addressing scheme used throughout SHRINCS.
\subsection{Hash Function}

SHRINCS uses SHA-256 truncated to \texttt{n} bytes for all internal hash operations. 
\begin{center}
\texttt{H(m) = SHA-256(m)[0:n]}
\end{center}

\subsection{Tweakable Hash Function}
To prevent multi-target attacks and ensure domain separation between the various components, we employ the tweakable hash function construction defined in SPHINCS+ \cite{sphincsplus_paper2019}:

\begin{center}
\texttt{Tw(PK.seed, ADRS, m) = H(PK.seed || ADRS || m)}
\end{center}

This usage of tweaked hashing is critical. It ensures that a hash computed at one position in the Merkle tree cannot be substituted for a hash at another position, effectively mitigating generic tree-grafting attacks.

\paragraph{Note:} For SHA-256, the compression function operates on 64-byte blocks. Since \texttt{PK.seed} is 16 bytes, we pad it to 64 bytes to enable precomputation:

\begin{center}
    \texttt{Tw(PK.seed, ADRS, m) = H(PK.seed || 0\string^48|| ADRS || m)}
\end{center}

Having the first 64 bytes constant (\texttt{PK.seed || 0\string^48}), we can precompute the compression function state once per key pair. Each tweak call then only needs to process the variable part (\texttt{ADRS || m}), reducing per-call overhead.

\subsection{Pseudorandom Functions}
SHRINCS uses two pseudorandom functions for different purposes: key derivation and getting the randomizer for the signature. 

For deterministic derivation of secret key material (WOTS+C chain starting values, FORS secret keys):

\begin{center}
    \texttt{PRF(PK.seed, ADRS, SK.seed) = Tw(PK.seed, ADRS, SK.seed)}
\end{center}

The address structure ADRS ensures that each derived value is unique and bound to its position in the scheme. 

For generating the randomizer \texttt{R} used in message hashing:

\begin{center}
    \texttt{PRF\_msg(SK.prf, PK.seed, opt\_rand, m) = HMAC-SHA-256(SK.prf, PK.seed || opt\_rand || m)[0:R\_SIZE]},
\end{center}
where \texttt{opt\_rand} is an optional randomness (\texttt{n} bytes), can be \texttt{PK.seed} for deterministic signing.
%\mknote{We need to unclude the counter in the PRF\_msg function for the stateless brunch, not in the message hash.}
\subsection{Message Hash Function}
SHRINCS uses different hashing strategies for WOTS+C and FORS+C, optimized for their respective use cases.

\subsubsection{WOTS+C: Two-Phase Hashing}
WOTS+C uses a two-phase approach for message hashing and grinding. This design is necessary because:
\begin{itemize}
    \item User messages need randomized hashing (phase 1 + phase 2)
    \item Internal XMSS signatures sign tree roots that are already hashes (phase 2 only)
\end{itemize}

\paragraph{Phase 1. Message Digest.} Computes a randomized hash of the user message:
\begin{center}
    \texttt{H\_msg(ADRS, R, PK.seed, PK.root, m) = SHA-256(ADRS || R || PK.seed || PK.root || m)[0:n]}
\end{center}
This is computed once per signature, producing an n-byte digest.

\paragraph{Phase 2. Grinding Hash.} Takes the digest and a counter:
\begin{center}
    \texttt{H\_grind(ADRS, digest, ctr) = SHA-256(ADRS || digest || ctr)[0:n]}
\end{center}
%\mknote{PKseed is needed as an input.}
This is computed repeatedly during grinding, but operates only on small fixed-size inputs (52 bytes for out parameters), making iterations fast.

\subsubsection{FORS+C: Single-Phase Hashing}
FORS+C always signs user messages (never tree roots), so it uses a simpler single-phase approach with the counter included directly:
\begin{center}
    \texttt{H\_msg\_fors(ADRS, R, ctr, PK.seed, PK.root, m) = SHA-256(ADRS || R || ctr || PK.seed || PK.root || m)[0:roundup((k*a)/8])}
\end{center}
%\mknote{We exlude the counter from here, we iterate the PRF\_msg with the counter, that samples R, and then we check if the sampled R works for the message hash.}

\subsection{Tweak Structure and Domain Separation}
Tweaks are structured as follows to ensure uniqueness across all scheme components:

\begin{center}
\texttt{ADRS = layer || tree\_address || type || words}
\end{center}
where: 
\begin{itemize}
    \item[] \texttt{layer}: (4 bytes) layer in hypertree (\texttt{0} = bottom, \texttt{d-1} = max)
    \item[] \texttt{tree\_address}: (12 bytes) tree address within layer
    \item[] \texttt{type}: (4 bytes) address type
    \item[] \texttt{words}: (12 bytes) type-dependent 12 bytes
\end{itemize}
A total \texttt{ADRS} size is 32 bytes.

\subsubsection{Address Types}
%\mknote{The grinding of FORS and FORS PRF becomes the same address type i think}
\begin{verbatim}
0x00: SF_WOTS_HASH: Stateful WOTS+C chain hashing
0x01: SF_WOTS_PK:   Stateful WOTS+C public key compression
0x02: SF_TREE:      Stateful tree internal nodes
0x03: SL_WOTS_HASH: Stateless WOTS+C chain hashing  
0x04: SL_WOTS_PK:   Stateless WOTS+C public key compression
0x05: SL_TREE:      Stateless tree internal nodes
0x06: FORS_HASH:    FORS leaf hashing
0x07: FORS_TREE:    FORS tree internal nodes
0x08: FORS_PK:      FORS roots compression
0x09: WOTS_GRIND:   WOTS+C grinding
0x0A: FORS_GRIND:   FORS+C grinding
0x0B: H_MSG:        Message hash
0x0C: WOTS_PRF:     PRF for WOTS+C key derivation
0x0D: FORS_PRF:     PRF for FORS secret key derivation
0x0E: ROOT:         Root public key computation
\end{verbatim}
Summarizing, \texttt{(0x00, 0x01, 0x02)} are stateful only types, \texttt{(0x03, 0x04, 0x05, 0x06, 0x07, 0x08, 0x0A)} are stateless only and \texttt{(0x09, 0x0B, 0x0C, 0x0D)} are shared types. 

\subsubsection{Address Manipulation Functions}
\begin{verbatim}
new_ADRS():
     return toByte(0, 32)

setLayerAddress(ADRS, layer):
    ADRS[0:4] ← toByte(layer, 4)

setTreeAddress(ADRS, tree_addr):
    ADRS[4:16] ← toByte(tree_addr, 12)

setTypeAndClear(ADRS, type):
    ADRS[16:20] ← toByte(type, 4)
    ADRS[20:32] ← toByte(0, 12)

setKeyPairAddress(ADRS, keypair):
    ADRS[20:24] ← toByte(keypair, 4)

setChainAddress(ADRS, chain):
    ADRS[24:28] ← toByte(chain, 4)

setHashAddress(ADRS, hash):
    ADRS[28:32] ← toByte(hash, 4)

setTreeHeight(ADRS, height):
    ADRS[24:28] ← toByte(height, 4)

setTreeIndex(ADRS, index):
    ADRS[28:32] ← toByte(index, 4)
\end{verbatim}

Note that \texttt{setTreeHeight} and \texttt{setChainAddress} both write to bytes \texttt{[24:28]}. These operations are context-dependent so there mustn't be a conflict.

\subsubsection{Type-Dependent Words}
The list of type-dependent words for address tweaking is the following:
\begin{table}[h!]
\centering
\begin{tabular}{c c c c c}
\hline
\textbf{Type} & \textbf{Name} & \textbf{word1 [20:24]} & \textbf{word2 [24:28]} & \textbf{word3 [28:32]} \\
\hline
0x00 & \texttt{SF\_WOTS\_HASH}  & keypair   & chain   & hash \\
0x01 & \texttt{SF\_WOTS\_PK}    & keypair   & 0       & 0 \\
0x02 & \texttt{SF\_TREE}        & keypair   & height  & index \\
0x03 & \texttt{SL\_WOTS\_HASH}  & keypair   & chain   & hash \\
0x04 & \texttt{SL\_WOTS\_PK}    & keypair   & 0       & 0 \\
0x05 & \texttt{SL\_TREE}        & keypair   & height  & index \\
0x06 & \texttt{FORS\_HASH}      & tree\_idx & 0       & leaf\_idx \\
0x07 & \texttt{FORS\_TREE}      & tree\_idx & height  & index \\
0x08 & \texttt{FORS\_PK}        & 0         & 0       & 0 \\
0x09 & \texttt{WOTS\_GRIND}     & keypair   & 0       & 0 \\
0x0A & \texttt{FORS\_GRIND}     & 0         & 0       & 0 \\
0x0B & \texttt{H\_MSG}          & 0         & 0       & 0 \\
0x0C & \texttt{WOTS\_PRF}       & keypair   & chain   & 0 \\
0x0D & \texttt{FORS\_PRF}       & tree\_idx & 0       & leaf\_idx \\
0x0E & \texttt{ROOT}            & 0         & 0       & 0 \\
\hline
\end{tabular}
\caption{Type-dependent interpretation of ADRS words \texttt{(bytes [20:32])}}
\end{table}